%% ma: L10-B02-0019
%% lop: 10
%% chuong: 1
%% bai: 2
%% dang: 10.2.9
%% loai: TLN
%% muc: TH
%% dap_an: "6"
%% nguon: "Ảnh giáo viên gửi 10/10/2026 (ThS. Nguyễn Hoàng Việt, Chương 2 Tập hợp), A-Đề 1, Phần 3 Bài 1"
%% kien_thuc: "Bài 2 (tập hợp và các phép toán trên tập hợp)"
%% trang_thai: cho-duyet
%% ly_do: "Dạng toán mới đề xuất 10.2.9: bài toán đếm số phần tử bằng biểu đồ Ven"
\begin{ex}
Lớp $10$A có $45$ học sinh trong đó có $12$ bạn học sinh giỏi môn Văn, $9$ bạn học sinh giỏi môn Địa lí, và $30$ bạn không giỏi môn học nào trong hai môn Văn, Địa lí. Hỏi lớp $10$A có bao nhiêu bạn học sinh giỏi cả hai môn Văn và Địa lí?
\shortans{6}
\loigiai{Số học sinh giỏi ít nhất một trong hai môn là $45-30=15$. Theo công thức $n(A\cup B)=n(A)+n(B)-n(A\cap B)$: $15=12+9-n(A\cap B)$, suy ra $n(A\cap B)=6$.}
\end{ex}
